\subsubsection{Subproceso 3.B: Confirmación y Registro Transaccional Oficial}

\begin{itemize}
    \item \textbf{Estudiante Regular:} Receptor final del Reporte de Matrícula, de la Boleta Oficial y de las acreditaciones correspondientes. Le corresponde verificar la exactitud de sus datos registrados, emplear la Boleta visada en los trámites que la exigen (becas, seguros, carné universitario) y canalizar ante su Facultad cualquier corrección necesaria, pues carece de edición directa sobre la matrícula ya asentada.

    \item \textbf{Oficina de Estadística de la Facultad:} Órgano de soporte operativo del cierre. Procesa los ajustes de última hora durante la semana de depuración, remite a la DIRCE el archivo digital de consolidación al concluir la semana, visa el Reporte y la Boleta de Matrícula, y formaliza ante la dependencia central la culminación del proceso de matrícula en la Facultad.

    \item \textbf{Dirección de Registro Central y Estadística (DIRCE):} Dependencia central del Vicerrectorado Académico que custodia la base de datos del SIGA. Efectúa el asentamiento transaccional definitivo, genera las boletas oficiales, consolida el padrón nacional de matriculados, pone a disposición de las Facultades el registro y la plataforma actualizada, y gestiona la interoperabilidad con las entidades estatales.

    \item \textbf{Docentes de Asignatura:} Receptores de los padrones oficiales de aula, insumos para el control de asistencias y el registro de calificaciones. No intervienen en el registro transaccional, pero la oportunidad de sus nóminas depende directamente del plazo de distribución del subproceso.

    \item \textbf{Organismos Nacionales (SUNEDU y MINEDU):} Entidades supervisoras y receptoras de la nómina consolidada de matrícula universitaria. Validan la información remitida para la estadística nacional de la educación superior y para los procesos de licenciamiento y carnetización.
\end{itemize}
